% ============ 6. Declaración de IA y reparto (S3) ============
\section{Declaración de uso de herramientas de IA y reparto del trabajo}

\subsection{Reparto del trabajo}

\begin{itemize}
  \item \textbf{Kevin Leonardo Chaparro Reyes} — rol principal: ingesta y
        auditoría de calidad de los datos (OE1); responsable de la
        reproducibilidad del pipeline; revisó el conjunto.
  \item \textbf{Valentina Muñoz Palma} — rol principal: análisis longitudinal
        (OE2: panel, transiciones y supervivencia) y programación estadística;
        revisó el conjunto.
  \item \textbf{Paula Margarita Triana Ancinez} — rol principal: estado del
        arte, análisis de equidad (OE3) y coordinación de la redacción y la
        sustentación; revisó el conjunto.
\end{itemize}
Todos los integrantes leyeron y pueden defender la totalidad del anteproyecto;
ninguna sección fue redactada por un solo integrante sin revisión de los demás.
\textit{(Ajustable a la división real del trabajo.)}

\subsection{Uso declarado de herramientas de IA}

Usamos herramientas de IA generativa y agentes autónomos como \textbf{apoyo},
no como autores. Específicamente:

\begin{itemize}
  \item \textbf{Cuánto se usó:} de forma \textbf{intensiva como asistente de
        análisis y redacción}: la mayoría de los borradores, fichas de
        auditoría y material de apoyo se generaron con IA; \textbf{todo el
        contenido final fue revisado, corregido y aprobado por los tres
        integrantes} antes de la entrega.
  \item \textbf{Para qué se usó:} análisis y documentación de código del
        repositorio de referencia, generación de fichas de auditoría, búsqueda
        y síntesis de literatura para el estado del arte, redacción de
        borradores y revisión de redacción, y propuesta de estructura del
        documento.
  \item \textbf{Validación en dos niveles, como lo solicitó el equipo:}
        (i) el \textbf{propio agente} que generó el material autoverificó cada
        afirmación \textbf{empíricamente} (ejecución real de los scripts, filas
        reales del XLSX, 133 commits con \texttt{git log}, exit codes de 16
        scripts) y con reglas automáticas; (ii) \textbf{agentes de validación
        independientes solicitados por el equipo} revisaron hechos, reglas de
        construcción, referencias (0 inventadas) y la rúbrica (bloques D y S);
        sus informes quedaron archivados en \texttt{data/validacion/}.
  \item \textbf{Qué se descartó porque estaba mal (casos concretos):}
        (1) un agente afirmó que el código comparaba la cadena \texttt{"SÍ"}
        con un dato \texttt{"SI"} y que la métrica de víctimas colapsaría a
        cero; al leer el código y ejecutarlo se verificó que la normalización
        previa lo resolvía: la afirmación se corrigió. (2) Un agente reportó un
        "orden alfabético" de archivos que en realidad era canónico por capa
        (artefacto de su propio verificador); se corrigió. (3) Un agente generó
        texto con un error de codificación (\texttt{SÃ} por \texttt{Sí}) que se
        detectó y corrigió en la revisión final.
\end{itemize}

\subsection{Riesgos éticos del proyecto y marco legal}

El proyecto trata \textbf{datos personales} del padrón público (identificadores
de persona, edad, geolocalización) e incluye \textbf{categorías sensibles}
auto-declaradas: grupo étnico, condición de discapacidad y condición de
víctima del conflicto. El marco legal aplicable en Colombia es la
\textbf{Ley 1581 de 2012} (protección de datos personales). Nuestras medidas:
(i) no publicar microdatos identificables: los análisis se publican agregados y
se mantiene la "gran tabla sin identificadores" como estándar, porque incluso
los datos seudonimizados pueden ser reidentificables en ciertas condiciones
(Rocher et al., 2019); (ii) declarar la finalidad y el tratamiento en la
documentación del proyecto; (iii) tratar las
variables sensibles con especial cuidado interpretativo (son auto-declaradas y
solo comparables en la ventana 2021). El riesgo de sesgo se mitiga declarando
explícitamente los supuestos y las limitaciones de cada indicador.
